\documentclass[11pt,a4paper]{article}
\usepackage[margin=2cm,top=1.8cm,bottom=1.8cm]{geometry}
\usepackage[utf8]{inputenc}
\usepackage[T1]{fontenc}
\usepackage{hyperref}
\usepackage{xcolor}
\usepackage{amssymb}

\title{Inverting the Formalization Workflow:\\
Prototyping an MPC Protocol in Rocq with an LLM Agent}
\author{Cheng-Hui Weng\\
Nagoya University}
\date{}

\begin{document}
\maketitle
\vspace{-1em}
\thispagestyle{empty}

\begin{abstract}
% OLD (trimmed for 2-page body)
% Formalization traditionally comes last, and this creates a high barrier when a
% construction depends on unfamiliar mathematics. I instead use formalization as
% the first step and treat it as a prototyping medium. In this workflow, an LLM
% agent drafts formal models, a theorem prover checks whether they stand on solid
% logical ground, and the human steers the process while learning from the output.
% I report on an ongoing Rocq formalization of SMC-PGG, a covering-space-based SMC
% protocol that spans five mathematical domains learned through formalization
% itself; the development currently includes 42 Theorem declarations. The paper
% closes with observations about human judgment, multi-domain lemma drafting by
% the LLM, and prover-backed certainty.
Formalization is usually the last step, but that ordering is costly when a
construction spans unfamiliar mathematics. I invert the order and use
formalization as a prototyping medium: an LLM drafts candidate models, Rocq
checks them, and the human steers revisions. The case study is an ongoing
formalization of SMC-PGG, a covering-space-based secure multi-party computation
framework with 42 Theorem declarations.
The main takeaway is methodological: prover feedback can guide
cross-domain learning and design decisions early, before full pen-and-paper
domain mastery.
\end{abstract}

\noindent\textbf{Keywords.} LLM-assisted formalization, Rocq, interactive theorem proving, secure multi-party computation, covering spaces.

\noindent\textbf{Introduction.}
% OLD (trimmed for 2-page body)
Formalization traditionally comes last. A researcher studies the relevant
mathematics, designs a construction informally, and only then translates the
result into a proof assistant as a final verification step. This ordering works
well when the mathematical landscape is already familiar, but it becomes a
serious obstacle when the construction draws on several domains that the
researcher has not yet studied. In this work, I reverse the order and use
formalization as a prototyping medium. An LLM agent\footnote{All LLM-assisted
work reported here used Claude (Anthropic), specifically Claude Opus 4.6.}
drafts formal models from natural-language direction, and a theorem prover
checks whether each draft stands on solid logical ground. In this setting, the
human learns from the structured output of each iteration. Definitions accepted
by the theorem prover reveal useful abstractions, and failed proof obligations can
signal missing assumptions, side conditions, or lemmas.

%%Formalization is often treated as end-stage verification. I use it earlier as a
%%design filter for ideas that span unfamiliar domains. An LLM
%%agent\footnote{All LLM-assisted work reported here used Claude (Anthropic),
%%specifically Claude Opus 4.6.} proposes formal candidates, Rocq validates or
%%rejects them, and each failure pinpoints missing assumptions or lemmas.
%
% This extended abstract makes three case-study contributions. First, it presents
% an inverted workflow in which formalization precedes domain mastery, following
% the pattern idea $\to$ formalize $\to$ learn, so that unfamiliar
% mathematical ideas can be tested early on solid logical ground. Second, it
% reports on an ongoing Rocq formalization of SMC-PGG, a secure multi-party computation
% protocol inspired by the covering space from topology;
% the development currently includes 42 Theorem declarations. Third, it
% draws methodological observations about the distinct roles of human judgment,
% multi-domain lemma drafting by the LLM, and prover-backed certainty.


This case study contributes an inverted workflow (``idea $\to$ formalize $\to$
learn``) and an ongoing Rocq formalization of SMC-PGG with 42 Theorem
declarations, highlighting how human judgment, LLM drafting, and prover checks
play distinct roles.

\noindent\textbf{Inverted Workflow.}

The starting point was a hypothesis from separate ongoing work on generalizing
secure multi-party computation (SMC). The hypothesis was that any surjective
relation can serve as a secret-sharing mechanism if legitimate parties possess a
trapdoor for reconstruction. Covering spaces in algebraic topology fit this
pattern. An index in the fiber over a basepoint can represent the secret, the
monodromy action can represent the computation, and a trapdoor inverting the
covering map can represent reconstruction. However, even formulating such
connections informally already requires time to understand each contributing
domain, and developing them to the level of a publishable pen-and-paper
argument takes much longer. Rather than waiting to master group theory,
algebraic geometry, and coding theory in advance, I used formalization in
Rocq/MathComp to test whether this connection makes sense on solid logical
ground.

From that point onward, formalization served not as a final verification step
but as an exploratory instrument. Rather than waiting for complete prior
mastery of every contributing domain, I used proof development to test whether
the core hypothesis could be realized as a concrete protocol with checkable
correctness and security guarantees.
% %
With this framing, I first asked the agent to assess the mathematical
plausibility of the secret-hiding assumption and to compile working notes on
implementations for sharing, computation, and reconstruction. Using those
notes, I then asked the agent to draft formalization plans and evaluate their
feasibility against available Rocq infrastructure, especially
MathComp~\cite{MathComp} and Infotheo~\cite{InfoTheo}. In parallel, the search
expanded into spectral graph theory, algebraic geometry, and coding theory
while I learned unfamiliar material.
% %
This process also exposed infrastructure boundaries that required explicit
axiomatization. For instance, the Riemann--Hurwitz formula is central to
linking group size with reconstruction constraints, yet a full mechanization
would itself be comparable to a separate research project. When a direction
appeared unlikely to meet the target guarantees, I required the agent to
produce structured tabular retrospectives that explicitly listed obstacles and
alternatives, preventing vague narrative evasion and forcing detailed
justification of the underlying reasoning.
% %
Formalization thus functioned as a methodological filter, separating components
ready for mechanization from those that should remain explicit assumptions at
this stage. The next section shows how this filter shaped the concrete SMC-PGG
case-study formalization.

%%The starting hypothesis was that surjective relations with a reconstruction
%%trapdoor can instantiate secret sharing. Covering spaces provide one candidate:
%%fiber indices encode secrets, monodromy models computation, and inversion of the
%%covering map models reconstruction. Instead of waiting for full prior mastery of
%%all background domains, I used Rocq/MathComp to test this structure directly.

%%The loop was simple: ask the agent for formal candidates, validate feasibility
%%against MathComp~\cite{MathComp} and Infotheo~\cite{InfoTheo}, and keep only
%%components that survived proof obligations. The process also exposed boundaries
%%that remain explicit assumptions (for example, results whose full mechanization
%%would be a separate project). In this sense, formalization acted as a
%%methodological filter before committing to a full protocol architecture.
%

\noindent\textbf{Case Study.}
To test the inverted workflow, I used it to build SMC-PGG (Secure Multi-party
Computation via Parametric Geometry Group), starting from the core idea that a
covering-space model can hide a secret as an index in the fiber over a base
point and evolve it through monodromy, formalized as a representation \(\rho : G \to S_N\).
Through a human--AI hybrid formalization loop, this idea grew into a parametric
mathematical framework before any venue- or domain-specific targeting;
AI-assisted literature search then revealed that the framework aligns naturally
with card-based SMC protocol research.
%
% OLD (trimmed for 2-page body)
% Concretely, the dealer distributes the secret as \(T\) shares placed at
% positions drawn from the sheet set \(\mathbb{I}_N\); each protocol step applies a
% permutation \(\rho(g)\) of \(\mathbb{I}_N\) that reshuffles these share positions while
% the reconstructed secret value stays invariant. This gives one parametric
% semantics for both single-round and repeated shuffles.
%
% The formalization organizes guarantees in three layers. The security layer
% tracks how far the induced shuffle distribution is from uniform and provides
% bounds on adversary advantage. The reconstruction layer captures classical
% \((k,T)\)-threshold behavior: any \(T\) shares reconstruct, while fewer than
% \(T\) reveal only bounded information. The combined layer connects security
% and threshold constraints through group-size and geometric conditions. Within
% this framework, I prove bounds on adversary search space, indistinguishable
% computation histories, privacy loss, and threshold margin, including the
% collusion bound
Concretely, the dealer places \(T\) shares (total share count) on sheet
positions \(\mathbb{I}_N\), and each protocol step applies \(\rho(g)\in S_N\),
reshuffling positions while preserving reconstruction semantics. The
formalization combines security and \((K,T)\)-threshold constraints (\(K\):
number of shares needed to reconstruct the secret) in one parametric model and
proves privacy/search-space bounds.

%, including
%
%\[
%\mathrm{TV} \le \varepsilon + \frac{2(T-1)}{N}.
%\]

%
The framework is validated on four instances: uniform shuffling
\cite{denBoer1989}, biased shuffling \cite{KimCetinkaya2025}, a classical
\(S_5\) construction, and an \(S_5 \times S_5\) construction. The first three fit the
simpler five-sheet setting; the \(S_5 \times S_5\) case crosses the group size/geometry
boundary and forces positive dropout tolerance \(T-K>0\). This provides evidence
that the central trade-off is realized on both sides of the boundary.
Except those four instances, I also show an impossibility
result: in the regular abelian setting, symmetry alone cannot deliver
non-trivial privacy/security gain.
%
Methodologically, the section also illustrates learning through formalization:
stronger arguments were introduced when simpler proof paths failed. The current
development contains 42 Theorem declarations across group theory, combinatorics,
algebraic geometry, coding theory, and information theory.

%%I also obtain a quantum-search scaling interpretation and an impossibility
%%result: in regular abelian settings, symmetry alone does not yield non-trivial
%%privacy gain. The framework is validated on four instances (including biased
%%shuffling and a product-group case that forces \(T-K>0\)), and the current
%%development contains 42 Theorem declarations.

\noindent\textbf{Observations.}
This case study suggests a methodological distinction from current lines of
automated-reasoning work that prioritize theorem-level automation, such as
whole-proof generation, informal-to-formal translation, or benchmark solving
\cite{first2023baldur,jiang2023dsp,alphaproof2024}. The workflow reported here
instead semi-automates framework construction: it supports exploration across
domains, iterative lemma drafting, and composition of many proved components
into one coherent protocol architecture. In this setting, the main challenge is
not only closing isolated goals, but also deciding which abstractions should
exist and how security, reconstruction, and algebraic constraints should fit
together.

The project also surfaced a practical architecture lesson. During repeated
pivots, some formally verified components became obsolete, over-general, or
temporarily disconnected from the main proof path. Examples include a
Cartier-Foata trace-counting theorem, a right-angled Artin group word-reordering
module, Massey's secret sharing from linear codes, and a weight bound for codes
on hyperelliptic curves. These compile and are available for future instances,
but do not load-bear in the current security or reconstruction proofs. This is
less a failure of proving than a normal software-architecture effect under
discovery-driven research: the framework requires continuous housekeeping to
deprecate stale artifacts, refactor interfaces, and keep dependencies and
theorem roles explicit, so the final development remains clear, auditable, and
internally consistent. 

%%Relative to theorem-level automation, this workflow targets framework
%%construction: selecting abstractions, drafting cross-domain lemmas, and
%%assembling a coherent security/reconstruction architecture. A practical lesson
%%is that discovery-driven formalization accumulates stale but compilable artifacts,
%%so continuous curation (deprecation, refactoring, explicit theorem roles) is
%%part of the method, not an afterthought.

\noindent\textbf{Related Work.}
% OLD (trimmed for 2-page body)
% Related work on LLM-assisted proving includes whole-proof generation
% \cite{first2023baldur}, informal-to-formal translation \cite{jiang2023dsp}, and
% retrieval-augmented search \cite{yang2023leandojo}. It also includes IDE
% integration \cite{song2024leancopilot} and reinforcement-learning-based training
% \cite{alphaproof2024}. Here, by contrast, the LLM primarily serves as a
% mechanism for domain-knowledge transfer in a human-directed exploratory
% workflow.
LLM-assisted proving work emphasizes proof generation, translation, retrieval,
and training \cite{first2023baldur,jiang2023dsp,yang2023leandojo,song2024leancopilot,alphaproof2024}.
In contrast, this project uses the LLM mainly for domain-knowledge transfer and
design-space exploration inside a prover-checked loop.

\noindent\textbf{Conclusion.}
% OLD (trimmed for 2-page body)
% This case study suggests a methodological claim about AI-assisted formalization.
% An LLM agent can be useful not only for proof search, but also for cross-domain
% exploration in an ongoing research project. In that setting, formalization need
% not be the final verification step. It can instead serve as a prototyping medium
% that tests whether a mathematical idea survives contact with a theorem prover
% while simultaneously teaching the surrounding domain. The traditional order of
% learn $\to$ design $\to$ formalize can therefore be inverted into idea
% $\to$ formalize $\to$ learn.
The case study supports a simple claim: AI-assisted formalization can guide
early-stage protocol design, not only late-stage proof search. In this mode,
formalization functions as a prototyping filter that both validates and teaches,
making ``idea $\to$ formalize $\to$ learn'' a practical research workflow.
From this perspective, the next step is twofold: to mature the current prototype
into a publication-ready development, and to test the same workflow on new ideas
in other domains.

\begingroup\footnotesize
\begin{thebibliography}{9}
\setlength{\itemsep}{0pt}\setlength{\parsep}{0pt}\setlength{\parskip}{0pt}

\bibitem{first2023baldur}
E.~First, M.~N.~Rabe, T.~Ringer, and Y.~Brun. Baldur: Whole-proof generation and
repair with large language models. In \textit{ESEC/FSE 2023}, 2023.

\bibitem{jiang2023dsp}
A.~Jiang et al. Draft, sketch, and prove: Guiding formal theorem provers with
informal proofs. In \textit{ICLR 2023}, 2023.

\bibitem{yang2023leandojo}
K.~Yang et al. LeanDojo: Theorem proving with retrieval-augmented language
models. In \textit{NeurIPS 2023 (Datasets and Benchmarks)}, 2023.

\bibitem{song2024leancopilot}
P.~Song, K.~Yang, and A.~Anandkumar. Lean Copilot: LLMs as copilots for theorem
proving in Lean.
\textit{arXiv:2404.12534}, 2024.

\bibitem{alphaproof2024}
Google DeepMind. AlphaProof and AlphaGeometry~2. Technical report, 2024.

\bibitem{MathComp}
The Mathematical Components Team. Mathematical Components, 2024.
\url{https://math-comp.github.io}

\bibitem{InfoTheo}
R.~Affeldt et al. A Coq Formalization of Information Theory and Linear
Error-Correcting Codes, 2024.
\url{https://github.com/affeldt-aist/infotheo}

\bibitem{barthe2011easycrypt}
G.~Barthe et al. Computer-aided security proofs for the working cryptographer.
In \textit{CRYPTO 2011}, LNCS 6841, Springer, 2011.

\bibitem{KimCetinkaya2025}
D.~H.~Kim and A.~Cetinkaya. Confidentiality in a Card-Based Protocol Under
Repeated Biased Shuffles.
\textit{arXiv:2511.05111 [cs.CR]}, 2025.

\bibitem{denBoer1989}
B.~den Boer. More Efficient Match-Making and Satisfiability: The Five Card
Trick. In \textit{Advances in Cryptology --- EUROCRYPT '89}, LNCS 434, pages
208--217, Springer, 1990.

\end{thebibliography}
\endgroup

\end{document}
